\chapter{概览}

% === 源文件: index.md ===
\section{OpenClaw}



\begin{quote}
\textit{"去壳！去壳！"} — 大概是一只太空龙虾说的
\end{quote}


适用于任何操作系统的 AI 智能体 Gateway 网关，支持 WhatsApp、Telegram、Discord、iMessage 等。
发送消息，随时随地获取智能体响应。通过插件可添加 Mattermost 等更多渠道。

安装 OpenClaw 并在几分钟内启动 Gateway 网关。
通过 \texttt{openclaw onboard} 和配对流程进行引导式设置。
启动浏览器仪表板，管理聊天、配置和会话。

OpenClaw 通过单个 Gateway 网关进程将聊天应用连接到 Pi 等编程智能体。它为 OpenClaw 助手提供支持，并支持本地或远程部署。

\subsection{工作原理}


\begin{verbatim}
% Mermaid 图表
flowchart LR
  A["Chat apps + plugins"] --> B["Gateway"]
  B --> C["Pi agent"]
  B --> D["CLI"]
  B --> E["Web Control UI"]
  B --> F["macOS app"]
  B --> G["iOS and Android nodes"]
\end{verbatim}


Gateway 网关是会话、路由和渠道连接的唯一事实来源。

\subsection{核心功能}


通过单个 Gateway 网关进程连接 WhatsApp、Telegram、Discord 和 iMessage。
通过扩展包添加 Mattermost 等更多渠道。
按智能体、工作区或发送者隔离会话。
发送和接收图片、音频和文档。
浏览器仪表板，用于聊天、配置、会话和节点管理。
配对 iOS 和 Android 节点，支持 Canvas。

\subsection{快速开始}


\begin{lstlisting}[language=bash]
    npm install -g openclaw@latest
\end{lstlisting}

\begin{lstlisting}[language=bash]
    openclaw onboard --install-daemon
\end{lstlisting}

\begin{lstlisting}[language=bash]
    openclaw channels login
    openclaw gateway --port 18789
\end{lstlisting}


需要完整的安装和开发环境设置？请参阅\href{/start/quickstart}{快速开始}。

\subsection{仪表板}


Gateway 网关启动后，打开浏览器控制界面。

\begin{itemize}
  \item 本地默认地址：http://127.0.0.1:18789/
  \item 远程访问：\href{/web}{Web 界面}和 \href{/gateway/tailscale}{Tailscale}
\end{itemize}



\subsection{配置（可选）}


配置文件位于 \texttt{\textasciitilde{}/.openclaw/openclaw.json}。

\begin{itemize}
  \item 如果你\textbf{不做任何修改}，OpenClaw 将使用内置的 Pi 二进制文件以 RPC 模式运行，并按发送者创建独立会话。
  \item 如果你想要限制访问，可以从 \texttt{channels.whatsapp.allowFrom} 和（针对群组的）提及规则开始配置。
\end{itemize}


示例：

\begin{lstlisting}[language=json]
{
  channels: {
    whatsapp: {
      allowFrom: ["+15555550123"],
      groups: { "*": { requireMention: true } },
    },
  },
  messages: { groupChat: { mentionPatterns: ["@openclaw"] } },
}
\end{lstlisting}


\subsection{从这里开始}


所有文档和指南，按用例分类。
核心 Gateway 网关设置、令牌和提供商配置。
SSH 和 tailnet 访问模式。
WhatsApp、Telegram、Discord 等渠道的具体设置。
iOS 和 Android 节点的配对与 Canvas 功能。
常见修复方法和故障排除入口。

\subsection{了解更多}


全部渠道、路由和媒体功能。
工作区隔离和按智能体的会话管理。
令牌、白名单和安全控制。
Gateway 网关诊断和常见错误。
项目起源、贡献者和许可证。


\clearpage

% === 源文件: brave-search.md ===
\section{Brave Search API}


OpenClaw 使用 Brave Search 作为 \texttt{web\_search} 的默认提供商。

\subsection{获取 API 密钥}


\begin{enumerate}
  \item 在 https://brave.com/search/api/ 创建 Brave Search API 账户
  \item 在控制面板中，选择 \textbf{Data for Search} 套餐并生成 API 密钥。
  \item 将密钥存储在配置中（推荐），或在 Gateway 网关环境中设置 \texttt{BRAVE\_API\_KEY}。
\end{enumerate}


\subsection{配置示例}


\begin{lstlisting}[language=json]
{
  tools: {
    web: {
      search: {
        provider: "brave",
        apiKey: "BRAVE_API_KEY_HERE",
        maxResults: 5,
        timeoutSeconds: 30,
      },
    },
  },
}
\end{lstlisting}


\subsection{注意事项}


\begin{itemize}
  \item Data for AI 套餐与 \texttt{web\_search} \textbf{不}兼容。
  \item Brave 提供免费层级和付费套餐；请查看 Brave API 门户了解当前限制。
\end{itemize}


请参阅 \href{/tools/web}{Web 工具} 了解完整的 web\_search 配置。


\clearpage

% === 源文件: date-time.md ===
\section{日期与时间}


OpenClaw 默认使用\textbf{主机本地时间作为传输时间戳}，并且\textbf{仅在系统提示词中使用用户时区}。
提供商时间戳会被保留，因此工具保持其原生语义（当前时间可通过 \texttt{session\_status} 获取）。

\subsection{消息信封（默认为本地时间）}


入站消息会附带一个时间戳（分钟精度）：

\begin{lstlisting}
[Provider ... 2026-01-05 16:26 PST] message text
\end{lstlisting}


此信封时间戳\textbf{默认为主机本地时间}，与提供商时区无关。

你可以覆盖此行为：

\begin{lstlisting}[language=json]
{
  agents: {
    defaults: {
      envelopeTimezone: "local", // "utc" | "local" | "user" | IANA 时区
      envelopeTimestamp: "on", // "on" | "off"
      envelopeElapsed: "on", // "on" | "off"
    },
  },
}
\end{lstlisting}


\begin{itemize}
  \item \texttt{envelopeTimezone: "utc"} 使用 UTC。
  \item \texttt{envelopeTimezone: "local"} 使用主机时区。
  \item \texttt{envelopeTimezone: "user"} 使用 \texttt{agents.defaults.userTimezone}（回退到主机时区）。
  \item 使用显式 IANA 时区（例如 \texttt{"America/Chicago"}）指定固定时区。
  \item \texttt{envelopeTimestamp: "off"} 从信封头中移除绝对时间戳。
  \item \texttt{envelopeElapsed: "off"} 移除已用时间后缀（\texttt{+2m} 样式）。
\end{itemize}


\subsubsection{示例}


\textbf{本地时间（默认）：}

\begin{lstlisting}
[WhatsApp +1555 2026-01-18 00:19 PST] hello
\end{lstlisting}


\textbf{用户时区：}

\begin{lstlisting}
[WhatsApp +1555 2026-01-18 00:19 CST] hello
\end{lstlisting}


\textbf{启用已用时间：}

\begin{lstlisting}
[WhatsApp +1555 +30s 2026-01-18T05:19Z] follow-up
\end{lstlisting}


\subsection{系统提示词：当前日期与时间}


如果已知用户时区，系统提示词会包含一个专门的\textbf{当前日期与时间}部分，其中仅包含\textbf{时区}（不含时钟/时间格式），以保持提示词缓存的稳定性：

\begin{lstlisting}
Time zone: America/Chicago
\end{lstlisting}


当智能体需要获取当前时间时，请使用 \texttt{session\_status} 工具；状态卡中包含时间戳行。

\subsection{系统事件行（默认为本地时间）}


插入到智能体上下文中的排队系统事件会带有时间戳前缀，使用与消息信封相同的时区选择（默认：主机本地时间）。

\begin{lstlisting}
System: [2026-01-12 12:19:17 PST] Model switched.
\end{lstlisting}


\subsubsection{配置用户时区和格式}


\begin{lstlisting}[language=json]
{
  agents: {
    defaults: {
      userTimezone: "America/Chicago",
      timeFormat: "auto", // auto | 12 | 24
    },
  },
}
\end{lstlisting}


\begin{itemize}
  \item \texttt{userTimezone} 设置提示词上下文中的\textbf{用户本地时区}。
  \item \texttt{timeFormat} 控制提示词中的 \textbf{12 小时/24 小时显示格式}。\texttt{auto} 跟随操作系统偏好设置。
\end{itemize}


\subsection{时间格式检测（auto）}


当 \texttt{timeFormat: "auto"} 时，OpenClaw 会检查操作系统偏好设置（macOS/Windows），并回退到区域格式。检测到的值会\textbf{按进程缓存}，以避免重复的系统调用。

\subsection{工具载荷 + 连接器（原始提供商时间 + 标准化字段）}


渠道工具返回\textbf{提供商原生时间戳}，并添加标准化字段以保持一致性：

\begin{itemize}
  \item \texttt{timestampMs}：纪元毫秒数（UTC）
  \item \texttt{timestampUtc}：ISO 8601 UTC 字符串
\end{itemize}


原始提供商字段会被保留，不会丢失任何数据。

\begin{itemize}
  \item Slack：来自 API 的类纪元字符串
  \item Discord：UTC ISO 时间戳
  \item Telegram/WhatsApp：提供商特定的数字/ISO 时间戳
\end{itemize}


如果需要本地时间，请使用已知时区在下游进行转换。

\subsection{相关文档}


\begin{itemize}
  \item \href{/concepts/system-prompt}{系统提示词}
  \item \href{/concepts/timezone}{时区}
  \item \href{/concepts/messages}{消息}
\end{itemize}



\clearpage

% === 源文件: logging.md ===
\section{日志}


OpenClaw 在两个地方记录日志：

\begin{itemize}
  \item \textbf{文件日志}（JSON 行）由 Gateway 网关写入。
  \item \textbf{控制台输出}显示在终端和控制 UI 中。
\end{itemize}


本页说明日志存放位置、如何读取日志以及如何配置日志级别和格式。

\subsection{日志存放位置}


默认情况下，Gateway 网关在以下位置写入滚动日志文件：

\texttt{/tmp/openclaw/openclaw-YYYY-MM-DD.log}

日期使用 Gateway 网关主机的本地时区。

你可以在 \texttt{\textasciitilde{}/.openclaw/openclaw.json} 中覆盖此设置：

\begin{lstlisting}[language=json]
{
  "logging": {
    "file": "/path/to/openclaw.log"
  }
}
\end{lstlisting}


\subsection{如何读取日志}


\subsubsection{CLI：实时跟踪（推荐）}


使用 CLI 通过 RPC 跟踪 Gateway 网关日志文件：

\begin{lstlisting}[language=bash]
openclaw logs --follow
\end{lstlisting}


输出模式：

\begin{itemize}
  \item \textbf{TTY 会话}：美观、彩色、结构化的日志行。
  \item \textbf{非 TTY 会话}：纯文本。
  \item \texttt{--json}：行分隔的 JSON（每行一个日志事件）。
  \item \texttt{--plain}：在 TTY 会话中强制纯文本。
  \item \texttt{--no-color}：禁用 ANSI 颜色。
\end{itemize}


在 JSON 模式下，CLI 输出带 \texttt{type} 标签的对象：

\begin{itemize}
  \item \texttt{meta}：流元数据（文件、游标、大小）
  \item \texttt{log}：解析的日志条目
  \item \texttt{notice}：截断/轮转提示
  \item \texttt{raw}：未解析的日志行
\end{itemize}


如果 Gateway 网关无法访问，CLI 会打印一个简短提示运行：

\begin{lstlisting}[language=bash]
openclaw doctor
\end{lstlisting}


\subsubsection{控制 UI（Web）}


控制 UI 的\textbf{日志}标签页使用 \texttt{logs.tail} 跟踪相同的文件。
参见 \href{/web/control-ui}{/web/control-ui} 了解如何打开它。

\subsubsection{仅渠道日志}


要过滤渠道活动（WhatsApp/Telegram 等），使用：

\begin{lstlisting}[language=bash]
openclaw channels logs --channel whatsapp
\end{lstlisting}


\subsection{日志格式}


\subsubsection{文件日志（JSONL）}


日志文件中的每一行都是一个 JSON 对象。CLI 和控制 UI 解析这些条目以渲染结构化输出（时间、级别、子系统、消息）。

\subsubsection{控制台输出}


控制台日志\textbf{感知 TTY}并格式化以提高可读性：

\begin{itemize}
  \item 子系统前缀（例如 \texttt{gateway/channels/whatsapp}）
  \item 级别着色（info/warn/error）
  \item 可选的紧凑或 JSON 模式
\end{itemize}


控制台格式由 \texttt{logging.consoleStyle} 控制。

\subsection{配置日志}


所有日志配置都在 \texttt{\textasciitilde{}/.openclaw/openclaw.json} 的 \texttt{logging} 下。

\begin{lstlisting}[language=json]
{
  "logging": {
    "level": "info",
    "file": "/tmp/openclaw/openclaw-YYYY-MM-DD.log",
    "consoleLevel": "info",
    "consoleStyle": "pretty",
    "redactSensitive": "tools",
    "redactPatterns": ["sk-.*"]
  }
}
\end{lstlisting}


\subsubsection{日志级别}


\begin{itemize}
  \item \texttt{logging.level}：\textbf{文件日志}（JSONL）级别。
  \item \texttt{logging.consoleLevel}：\textbf{控制台}详细程度级别。
\end{itemize}


\texttt{--verbose} 仅影响控制台输出；它不改变文件日志级别。

\subsubsection{控制台样式}


\texttt{logging.consoleStyle}：

\begin{itemize}
  \item \texttt{pretty}：人类友好、彩色、带时间戳。
  \item \texttt{compact}：更紧凑的输出（最适合长会话）。
  \item \texttt{json}：每行 JSON（用于日志处理器）。
\end{itemize}


\subsubsection{脱敏}


工具摘要可以在敏感令牌输出到控制台之前对其进行脱敏：

\begin{itemize}
  \item \texttt{logging.redactSensitive}：\texttt{off} | \texttt{tools}（默认：\texttt{tools}）
  \item \texttt{logging.redactPatterns}：用于覆盖默认集的正则表达式字符串列表
\end{itemize}


脱敏仅影响\textbf{控制台输出}，不会改变文件日志。

\subsection{诊断 + OpenTelemetry}


诊断是用于模型运行\textbf{和}消息流遥测（webhooks、队列、会话状态）的结构化、机器可读事件。它们\textbf{不}替代日志；它们存在是为了向指标、追踪和其他导出器提供数据。

诊断事件在进程内发出，但导出器仅在启用诊断 + 导出器插件时才附加。

\subsubsection{OpenTelemetry 与 OTLP}


\begin{itemize}
  \item \textbf{OpenTelemetry（OTel）}：追踪、指标和日志的数据模型 + SDK。
  \item \textbf{OTLP}：用于将 OTel 数据导出到收集器/后端的线路协议。
  \item OpenClaw 目前通过 \textbf{OTLP/HTTP（protobuf）} 导出。
\end{itemize}


\subsubsection{导出的信号}


\begin{itemize}
  \item \textbf{指标}：计数器 + 直方图（令牌使用、消息流、队列）。
  \item \textbf{追踪}：模型使用 + webhook/消息处理的 span。
  \item \textbf{日志}：启用 \texttt{diagnostics.otel.logs} 时通过 OTLP 导出。日志量可能很大；请注意 \texttt{logging.level} 和导出器过滤器。
\end{itemize}


\subsubsection{诊断事件目录}


模型使用：

\begin{itemize}
  \item \texttt{model.usage}：令牌、成本、持续时间、上下文、提供商/模型/渠道、会话 ID。
\end{itemize}


消息流：

\begin{itemize}
  \item \texttt{webhook.received}：每渠道的 webhook 入口。
  \item \texttt{webhook.processed}：webhook 已处理 + 持续时间。
  \item \texttt{webhook.error}：webhook 处理程序错误。
  \item \texttt{message.queued}：消息入队等待处理。
  \item \texttt{message.processed}：结果 + 持续时间 + 可选错误。
\end{itemize}


队列 + 会话：

\begin{itemize}
  \item \texttt{queue.lane.enqueue}：命令队列通道入队 + 深度。
  \item \texttt{queue.lane.dequeue}：命令队列通道出队 + 等待时间。
  \item \texttt{session.state}：会话状态转换 + 原因。
  \item \texttt{session.stuck}：会话卡住警告 + 持续时间。
  \item \texttt{run.attempt}：运行重试/尝试元数据。
  \item \texttt{diagnostic.heartbeat}：聚合计数器（webhooks/队列/会话）。
\end{itemize}


\subsubsection{启用诊断（无导出器）}


如果你希望诊断事件可用于插件或自定义接收器，使用此配置：

\begin{lstlisting}[language=json]
{
  "diagnostics": {
    "enabled": true
  }
}
\end{lstlisting}


\subsubsection{诊断标志（定向日志）}


使用标志在不提高 \texttt{logging.level} 的情况下开启额外的定向调试日志。
标志不区分大小写，支持通配符（例如 \texttt{telegram.\textit{} 或 \texttt{}}）。

\begin{lstlisting}[language=json]
{
  "diagnostics": {
    "flags": ["telegram.http"]
  }
}
\end{lstlisting}


环境变量覆盖（一次性）：

\begin{lstlisting}
OPENCLAW_DIAGNOSTICS=telegram.http,telegram.payload
\end{lstlisting}


注意：

\begin{itemize}
  \item 标志日志进入标准日志文件（与 \texttt{logging.file} 相同）。
  \item 输出仍根据 \texttt{logging.redactSensitive} 进行脱敏。
  \item 完整指南：\href{/diagnostics/flags}{/diagnostics/flags}。
\end{itemize}


\subsubsection{导出到 OpenTelemetry}


诊断可以通过 \texttt{diagnostics-otel} 插件（OTLP/HTTP）导出。这适用于任何接受 OTLP/HTTP 的 OpenTelemetry 收集器/后端。

\begin{lstlisting}[language=json]
{
  "plugins": {
    "allow": ["diagnostics-otel"],
    "entries": {
      "diagnostics-otel": {
        "enabled": true
      }
    }
  },
  "diagnostics": {
    "enabled": true,
    "otel": {
      "enabled": true,
      "endpoint": "http://otel-collector:4318",
      "protocol": "http/protobuf",
      "serviceName": "openclaw-gateway",
      "traces": true,
      "metrics": true,
      "logs": true,
      "sampleRate": 0.2,
      "flushIntervalMs": 60000
    }
  }
}
\end{lstlisting}


注意：

\begin{itemize}
  \item 你也可以使用 \texttt{openclaw plugins enable diagnostics-otel} 启用插件。
  \item \texttt{protocol} 目前仅支持 \texttt{http/protobuf}。\texttt{grpc} 被忽略。
  \item 指标包括令牌使用、成本、上下文大小、运行持续时间和消息流计数器/直方图（webhooks、队列、会话状态、队列深度/等待）。
  \item 追踪/指标可以通过 \texttt{traces} / \texttt{metrics} 切换（默认：开启）。启用时，追踪包括模型使用 span 加上 webhook/消息处理 span。
  \item 当你的收集器需要认证时设置 \texttt{headers}。
  \item 支持的环境变量：\texttt{OTEL\_EXPORTER\_OTLP\_ENDPOINT}、\texttt{OTEL\_SERVICE\_NAME}、\texttt{OTEL\_EXPORTER\_OTLP\_PROTOCOL}。
\end{itemize}


\subsubsection{导出的指标（名称 + 类型）}


模型使用：

\begin{itemize}
  \item \texttt{openclaw.tokens}（计数器，属性：\texttt{openclaw.token}、\texttt{openclaw.channel}、\texttt{openclaw.provider}、\texttt{openclaw.model}）
  \item \texttt{openclaw.cost.usd}（计数器，属性：\texttt{openclaw.channel}、\texttt{openclaw.provider}、\texttt{openclaw.model}）
  \item \texttt{openclaw.run.duration\_ms}（直方图，属性：\texttt{openclaw.channel}、\texttt{openclaw.provider}、\texttt{openclaw.model}）
  \item \texttt{openclaw.context.tokens}（直方图，属性：\texttt{openclaw.context}、\texttt{openclaw.channel}、\texttt{openclaw.provider}、\texttt{openclaw.model}）
\end{itemize}


消息流：

\begin{itemize}
  \item \texttt{openclaw.webhook.received}（计数器，属性：\texttt{openclaw.channel}、\texttt{openclaw.webhook}）
  \item \texttt{openclaw.webhook.error}（计数器，属性：\texttt{openclaw.channel}、\texttt{openclaw.webhook}）
  \item \texttt{openclaw.webhook.duration\_ms}（直方图，属性：\texttt{openclaw.channel}、\texttt{openclaw.webhook}）
  \item \texttt{openclaw.message.queued}（计数器，属性：\texttt{openclaw.channel}、\texttt{openclaw.source}）
  \item \texttt{openclaw.message.processed}（计数器，属性：\texttt{openclaw.channel}、\texttt{openclaw.outcome}）
  \item \texttt{openclaw.message.duration\_ms}（直方图，属性：\texttt{openclaw.channel}、\texttt{openclaw.outcome}）
\end{itemize}


队列 + 会话：

\begin{itemize}
  \item \texttt{openclaw.queue.lane.enqueue}（计数器，属性：\texttt{openclaw.lane}）
  \item \texttt{openclaw.queue.lane.dequeue}（计数器，属性：\texttt{openclaw.lane}）
  \item \texttt{openclaw.queue.depth}（直方图，属性：\texttt{openclaw.lane} 或 \texttt{openclaw.channel=heartbeat}）
  \item \texttt{openclaw.queue.wait\_ms}（直方图，属性：\texttt{openclaw.lane}）
  \item \texttt{openclaw.session.state}（计数器，属性：\texttt{openclaw.state}、\texttt{openclaw.reason}）
  \item \texttt{openclaw.session.stuck}（计数器，属性：\texttt{openclaw.state}）
  \item \texttt{openclaw.session.stuck\_age\_ms}（直方图，属性：\texttt{openclaw.state}）
  \item \texttt{openclaw.run.attempt}（计数器，属性：\texttt{openclaw.attempt}）
\end{itemize}


\subsubsection{导出的 span（名称 + 关键属性）}


\begin{itemize}
  \item \texttt{openclaw.model.usage}
  \item \texttt{openclaw.channel}、\texttt{openclaw.provider}、\texttt{openclaw.model}
  \item \texttt{openclaw.sessionKey}、\texttt{openclaw.sessionId}
  \item \texttt{openclaw.tokens.*}（input/output/cache\_read/cache\_write/total）
  \item \texttt{openclaw.webhook.processed}
  \item \texttt{openclaw.channel}、\texttt{openclaw.webhook}、\texttt{openclaw.chatId}
  \item \texttt{openclaw.webhook.error}
  \item \texttt{openclaw.channel}、\texttt{openclaw.webhook}、\texttt{openclaw.chatId}、\texttt{openclaw.error}
  \item \texttt{openclaw.message.processed}
  \item \texttt{openclaw.channel}、\texttt{openclaw.outcome}、\texttt{openclaw.chatId}、\texttt{openclaw.messageId}、\texttt{openclaw.sessionKey}、\texttt{openclaw.sessionId}、\texttt{openclaw.reason}
  \item \texttt{openclaw.session.stuck}
  \item \texttt{openclaw.state}、\texttt{openclaw.ageMs}、\texttt{openclaw.queueDepth}、\texttt{openclaw.sessionKey}、\texttt{openclaw.sessionId}
\end{itemize}


\subsubsection{采样 + 刷新}


\begin{itemize}
  \item 追踪采样：\texttt{diagnostics.otel.sampleRate}（0.0–1.0，仅根 span）。
  \item 指标导出间隔：\texttt{diagnostics.otel.flushIntervalMs}（最小 1000ms）。
\end{itemize}


\subsubsection{协议说明}


\begin{itemize}
  \item OTLP/HTTP 端点可以通过 \texttt{diagnostics.otel.endpoint} 或 \texttt{OTEL\_EXPORTER\_OTLP\_ENDPOINT} 设置。
  \item 如果端点已包含 \texttt{/v1/traces} 或 \texttt{/v1/metrics}，则按原样使用。
  \item 如果端点已包含 \texttt{/v1/logs}，则按原样用于日志。
  \item \texttt{diagnostics.otel.logs} 为主日志器输出启用 OTLP 日志导出。
\end{itemize}


\subsubsection{日志导出行为}


\begin{itemize}
  \item OTLP 日志使用与写入 \texttt{logging.file} 相同的结构化记录。
  \item 遵守 \texttt{logging.level}（文件日志级别）。控制台脱敏\textbf{不}适用于 OTLP 日志。
  \item 高流量安装应优先使用 OTLP 收集器采样/过滤。
\end{itemize}


\subsection{故障排除提示}


\begin{itemize}
  \item \textbf{Gateway 网关无法访问？} 先运行 \texttt{openclaw doctor}。
  \item \textbf{日志为空？} 检查 Gateway 网关是否正在运行并写入 \texttt{logging.file} 中的文件路径。
  \item \textbf{需要更多细节？} 将 \texttt{logging.level} 设置为 \texttt{debug} 或 \texttt{trace} 并重试。
\end{itemize}



\clearpage

% === 源文件: network.md ===
\section{网络中心}


本中心汇集了 OpenClaw 如何在 localhost、局域网和 tailnet 之间连接、配对和保护设备的核心文档。

\subsection{核心模型}


\begin{itemize}
  \item \href{/concepts/architecture}{Gateway 网关架构}
  \item \href{/gateway/protocol}{Gateway 网关协议}
  \item \href{/gateway}{Gateway 网关运维手册}
  \item \href{/web}{Web 接口 + 绑定模式}
\end{itemize}


\subsection{配对 + 身份}


\begin{itemize}
  \item \href{/channels/pairing}{配对概述（私信 + 节点）}
  \item \href{/gateway/pairing}{Gateway 网关拥有的节点配对}
  \item \href{/cli/devices}{Devices CLI（配对 + token 轮换）}
  \item \href{/cli/pairing}{Pairing CLI（私信审批）}
\end{itemize}


本地信任：

\begin{itemize}
  \item 本地连接（loopback 或 Gateway 网关主机自身的 tailnet 地址）可以自动批准配对，以保持同主机用户体验的流畅性。
  \item 非本地的 tailnet/局域网客户端仍需要显式的配对批准。
\end{itemize}


\subsection{设备发现 + 传输协议}


\begin{itemize}
  \item \href{/gateway/discovery}{设备发现与传输协议}
  \item \href{/gateway/bonjour}{Bonjour / mDNS}
  \item \href{/gateway/remote}{远程访问（SSH）}
  \item \href{/gateway/tailscale}{Tailscale}
\end{itemize}


\subsection{节点 + 传输协议}


\begin{itemize}
  \item \href{/nodes}{节点概述}
  \item \href{/gateway/bridge-protocol}{桥接协议（旧版节点）}
  \item \href{/platforms/ios}{节点运维手册：iOS}
  \item \href{/platforms/android}{节点运维手册：Android}
\end{itemize}


\subsection{安全}


\begin{itemize}
  \item \href{/gateway/security}{安全概述}
  \item \href{/gateway/configuration}{Gateway 网关配置参考}
  \item \href{/gateway/troubleshooting}{故障排除}
  \item \href{/gateway/doctor}{Doctor}
\end{itemize}



\clearpage

% === 源文件: perplexity.md ===
\section{Perplexity Sonar}


OpenClaw 可以使用 Perplexity Sonar 作为 \texttt{web\_search} 工具。你可以通过 Perplexity 的直连 API 或通过 OpenRouter 连接。

\subsection{API 选项}


\subsubsection{Perplexity（直连）}


\begin{itemize}
  \item Base URL：https://api.perplexity.ai
  \item 环境变量：\texttt{PERPLEXITY\_API\_KEY}
\end{itemize}


\subsubsection{OpenRouter（替代方案）}


\begin{itemize}
  \item Base URL：https://openrouter.ai/api/v1
  \item 环境变量：\texttt{OPENROUTER\_API\_KEY}
  \item 支持预付费/加密货币积分。
\end{itemize}


\subsection{配置示例}


\begin{lstlisting}[language=json]
{
  tools: {
    web: {
      search: {
        provider: "perplexity",
        perplexity: {
          apiKey: "pplx-...",
          baseUrl: "https://api.perplexity.ai",
          model: "perplexity/sonar-pro",
        },
      },
    },
  },
}
\end{lstlisting}


\subsection{从 Brave 切换}


\begin{lstlisting}[language=json]
{
  tools: {
    web: {
      search: {
        provider: "perplexity",
        perplexity: {
          apiKey: "pplx-...",
          baseUrl: "https://api.perplexity.ai",
        },
      },
    },
  },
}
\end{lstlisting}


如果同时设置了 \texttt{PERPLEXITY\_API\_KEY} 和 \texttt{OPENROUTER\_API\_KEY}，请设置 \texttt{tools.web.search.perplexity.baseUrl}（或 \texttt{tools.web.search.perplexity.apiKey}）以消除歧义。

如果未设置 base URL，OpenClaw 会根据 API 密钥来源选择默认值：

\begin{itemize}
  \item \texttt{PERPLEXITY\_API\_KEY} 或 \texttt{pplx-...} → 直连 Perplexity（\texttt{https://api.perplexity.ai}）
  \item \texttt{OPENROUTER\_API\_KEY} 或 \texttt{sk-or-...} → OpenRouter（\texttt{https://openrouter.ai/api/v1}）
  \item 未知密钥格式 → OpenRouter（安全回退）
\end{itemize}


\subsection{模型}


\begin{itemize}
  \item \texttt{perplexity/sonar} — 带网络搜索的快速问答
  \item \texttt{perplexity/sonar-pro}（默认） — 多步推理 + 网络搜索
  \item \texttt{perplexity/sonar-reasoning-pro} — 深度研究
\end{itemize}


请参阅 \href{/tools/web}{Web 工具} 了解 web\_search 配置详情。


\clearpage

% === 源文件: pi-dev.md ===
\section{Pi 开发工作流程}


本指南总结了在 OpenClaw 中开发 Pi 集成的合理工作流程。

\subsection{类型检查和代码检查}


\begin{itemize}
  \item 类型检查和构建：\texttt{pnpm build}
  \item 代码检查：\texttt{pnpm lint}
  \item 格式检查：\texttt{pnpm format}
  \item 推送前完整检查：\texttt{pnpm lint \&\& pnpm build \&\& pnpm test}
\end{itemize}


\subsection{运行 Pi 测试}


使用专用脚本运行 Pi 集成测试集：

\begin{lstlisting}[language=bash]
scripts/pi/run-tests.sh
\end{lstlisting}


要包含执行真实提供商行为的实时测试：

\begin{lstlisting}[language=bash]
scripts/pi/run-tests.sh --live
\end{lstlisting}


该脚本通过以下 glob 模式运行所有 Pi 相关的单元测试：

\begin{itemize}
  \item \texttt{src/agents/pi-*.test.ts}
  \item \texttt{src/agents/pi-embedded-*.test.ts}
  \item \texttt{src/agents/pi-tools*.test.ts}
  \item \texttt{src/agents/pi-settings.test.ts}
  \item \texttt{src/agents/pi-tool-definition-adapter.test.ts}
  \item \texttt{src/agents/pi-extensions/*.test.ts}
\end{itemize}


\subsection{手动测试}


推荐流程：

\begin{itemize}
  \item 以开发模式运行 Gateway 网关：
  \item \texttt{pnpm gateway:dev}
  \item 直接触发智能体：
  \item \texttt{pnpm openclaw agent --message "Hello" --thinking low}
  \item 使用 TUI 进行交互式调试：
  \item \texttt{pnpm tui}
\end{itemize}


对于工具调用行为，提示执行 \texttt{read} 或 \texttt{exec} 操作，以便查看工具流式传输和负载处理。

\subsection{完全重置}


状态存储在 OpenClaw 状态目录下。默认为 \texttt{\textasciitilde{}/.openclaw}。如果设置了 \texttt{OPENCLAW\_STATE\_DIR}，则使用该目录。

要重置所有内容：

\begin{itemize}
  \item \texttt{openclaw.json} 用于配置
  \item \texttt{credentials/} 用于认证配置文件和 token
  \item \texttt{agents//sessions/} 用于智能体会话历史
  \item \texttt{agents//sessions.json} 用于会话索引
  \item \texttt{sessions/} 如果存在旧版路径
  \item \texttt{workspace/} 如果你想要一个空白工作区
\end{itemize}


如果只想重置会话，删除该智能体的 \texttt{agents//sessions/} 和 \texttt{agents//sessions.json}。如果不想重新认证，保留 \texttt{credentials/}。

\subsection{参考资料}


\begin{itemize}
  \item https://docs.openclaw.ai/testing
  \item https://docs.openclaw.ai/start/getting-started
\end{itemize}



\clearpage

% === 源文件: pi.md ===
\section{Pi 集成架构}


本文档描述了 OpenClaw 如何与 \href{https://github.com/badlogic/pi-mono/tree/main/packages/coding-agent}{pi-coding-agent} 及其相关包（\texttt{pi-ai}、\texttt{pi-agent-core}、\texttt{pi-tui}）集成以实现其 AI 智能体能力。

\subsection{概述}


OpenClaw 使用 pi SDK 将 AI 编码智能体嵌入到其消息 Gateway 网关架构中。OpenClaw 不是将 pi 作为子进程生成或使用 RPC 模式，而是通过 \texttt{createAgentSession()} 直接导入并实例化 pi 的 \texttt{AgentSession}。这种嵌入式方法提供了：

\begin{itemize}
  \item 对会话生命周期和事件处理的完全控制
  \item 自定义工具注入（消息、沙箱、渠道特定操作）
  \item 每个渠道/上下文的系统提示自定义
  \item 支持分支/压缩的会话持久化
  \item 带故障转移的多账户认证配置文件轮换
  \item 与提供商无关的模型切换
\end{itemize}


\subsection{包依赖}


\begin{lstlisting}[language=json]
{
  "@mariozechner/pi-agent-core": "0.49.3",
  "@mariozechner/pi-ai": "0.49.3",
  "@mariozechner/pi-coding-agent": "0.49.3",
  "@mariozechner/pi-tui": "0.49.3"
}
\end{lstlisting}



\begin{table}[h]
\centering
\small
\begin{tabular}{|l|l|}
\hline
包 & 用途 \\
\hline
\texttt{pi-ai} & 核心 LLM 抽象：\texttt{Model}、\texttt{streamSimple}、消息类型、提供商 API \\
\hline
\texttt{pi-agent-core} & 智能体循环、工具执行、\texttt{AgentMessage} 类型 \\
\hline
\texttt{pi-coding-agent} & 高级 SDK：\texttt{createAgentSession}、\texttt{SessionManager}、\texttt{AuthStorage}、\texttt{ModelRegistry}、内置工具 \\
\hline
\texttt{pi-tui} & 终端 UI 组件（用于 OpenClaw 的本地 TUI 模式） \\
\hline
\end{tabular}
\end{table}



\subsection{文件结构}


\begin{lstlisting}
src/agents/
├── pi-embedded-runner.ts          # Re-exports from pi-embedded-runner/
├── pi-embedded-runner/
│   ├── run.ts                     # Main entry: runEmbeddedPiAgent()
│   ├── run/
│   │   ├── attempt.ts             # Single attempt logic with session setup
│   │   ├── params.ts              # RunEmbeddedPiAgentParams type
│   │   ├── payloads.ts            # Build response payloads from run results
│   │   ├── images.ts              # Vision model image injection
│   │   └── types.ts               # EmbeddedRunAttemptResult
│   ├── abort.ts                   # Abort error detection
│   ├── cache-ttl.ts               # Cache TTL tracking for context pruning
│   ├── compact.ts                 # Manual/auto compaction logic
│   ├── extensions.ts              # Load pi extensions for embedded runs
│   ├── extra-params.ts            # Provider-specific stream params
│   ├── google.ts                  # Google/Gemini turn ordering fixes
│   ├── history.ts                 # History limiting (DM vs group)
│   ├── lanes.ts                   # Session/global command lanes
│   ├── logger.ts                  # Subsystem logger
│   ├── model.ts                   # Model resolution via ModelRegistry
│   ├── runs.ts                    # Active run tracking, abort, queue
│   ├── sandbox-info.ts            # Sandbox info for system prompt
│   ├── session-manager-cache.ts   # SessionManager instance caching
│   ├── session-manager-init.ts    # Session file initialization
│   ├── system-prompt.ts           # System prompt builder
│   ├── tool-split.ts              # Split tools into builtIn vs custom
│   ├── types.ts                   # EmbeddedPiAgentMeta, EmbeddedPiRunResult
│   └── utils.ts                   # ThinkLevel mapping, error description
├── pi-embedded-subscribe.ts       # Session event subscription/dispatch
├── pi-embedded-subscribe.types.ts # SubscribeEmbeddedPiSessionParams
├── pi-embedded-subscribe.handlers.ts # Event handler factory
├── pi-embedded-subscribe.handlers.lifecycle.ts
├── pi-embedded-subscribe.handlers.types.ts
├── pi-embedded-block-chunker.ts   # Streaming block reply chunking
├── pi-embedded-messaging.ts       # Messaging tool sent tracking
├── pi-embedded-helpers.ts         # Error classification, turn validation
├── pi-embedded-helpers/           # Helper modules
├── pi-embedded-utils.ts           # Formatting utilities
├── pi-tools.ts                    # createOpenClawCodingTools()
├── pi-tools.abort.ts              # AbortSignal wrapping for tools
├── pi-tools.policy.ts             # Tool allowlist/denylist policy
├── pi-tools.read.ts               # Read tool customizations
├── pi-tools.schema.ts             # Tool schema normalization
├── pi-tools.types.ts              # AnyAgentTool type alias
├── pi-tool-definition-adapter.ts  # AgentTool -> ToolDefinition adapter
├── pi-settings.ts                 # Settings overrides
├── pi-extensions/                 # Custom pi extensions
│   ├── compaction-safeguard.ts    # Safeguard extension
│   ├── compaction-safeguard-runtime.ts
│   ├── context-pruning.ts         # Cache-TTL context pruning extension
│   └── context-pruning/
├── model-auth.ts                  # Auth profile resolution
├── auth-profiles.ts               # Profile store, cooldown, failover
├── model-selection.ts             # Default model resolution
├── models-config.ts               # models.json generation
├── model-catalog.ts               # Model catalog cache
├── context-window-guard.ts        # Context window validation
├── failover-error.ts              # FailoverError class
├── defaults.ts                    # DEFAULT_PROVIDER, DEFAULT_MODEL
├── system-prompt.ts               # buildAgentSystemPrompt()
├── system-prompt-params.ts        # System prompt parameter resolution
├── system-prompt-report.ts        # Debug report generation
├── tool-summaries.ts              # Tool description summaries
├── tool-policy.ts                 # Tool policy resolution
├── transcript-policy.ts           # Transcript validation policy
├── skills.ts                      # Skill snapshot/prompt building
├── skills/                        # Skill subsystem
├── sandbox.ts                     # Sandbox context resolution
├── sandbox/                       # Sandbox subsystem
├── channel-tools.ts               # Channel-specific tool injection
├── openclaw-tools.ts              # OpenClaw-specific tools
├── bash-tools.ts                  # exec/process tools
├── apply-patch.ts                 # apply_patch tool (OpenAI)
├── tools/                         # Individual tool implementations
│   ├── browser-tool.ts
│   ├── canvas-tool.ts
│   ├── cron-tool.ts
│   ├── discord-actions*.ts
│   ├── gateway-tool.ts
│   ├── image-tool.ts
│   ├── message-tool.ts
│   ├── nodes-tool.ts
│   ├── session*.ts
│   ├── slack-actions.ts
│   ├── telegram-actions.ts
│   ├── web-*.ts
│   └── whatsapp-actions.ts
└── ...
\end{lstlisting}


\subsection{核心集成流程}


\subsubsection{1. 运行嵌入式智能体}


主入口点是 \texttt{pi-embedded-runner/run.ts} 中的 \texttt{runEmbeddedPiAgent()}：

\begin{lstlisting}[language=typescript]
import { runEmbeddedPiAgent } from "./agents/pi-embedded-runner.js";

const result = await runEmbeddedPiAgent({
  sessionId: "user-123",
  sessionKey: "main:whatsapp:+1234567890",
  sessionFile: "/path/to/session.jsonl",
  workspaceDir: "/path/to/workspace",
  config: openclawConfig,
  prompt: "Hello, how are you?",
  provider: "anthropic",
  model: "claude-sonnet-4-20250514",
  timeoutMs: 120_000,
  runId: "run-abc",
  onBlockReply: async (payload) => {
    await sendToChannel(payload.text, payload.mediaUrls);
  },
});
\end{lstlisting}


\subsubsection{2. 会话创建}


在 \texttt{runEmbeddedAttempt()}（由 \texttt{runEmbeddedPiAgent()} 调用）内部，使用 pi SDK：

\begin{lstlisting}[language=typescript]
import {
  createAgentSession,
  DefaultResourceLoader,
  SessionManager,
  SettingsManager,
} from "@mariozechner/pi-coding-agent";

const resourceLoader = new DefaultResourceLoader({
  cwd: resolvedWorkspace,
  agentDir,
  settingsManager,
  additionalExtensionPaths,
});
await resourceLoader.reload();

const { session } = await createAgentSession({
  cwd: resolvedWorkspace,
  agentDir,
  authStorage: params.authStorage,
  modelRegistry: params.modelRegistry,
  model: params.model,
  thinkingLevel: mapThinkingLevel(params.thinkLevel),
  tools: builtInTools,
  customTools: allCustomTools,
  sessionManager,
  settingsManager,
  resourceLoader,
});

applySystemPromptOverrideToSession(session, systemPromptOverride);
\end{lstlisting}


\subsubsection{3. 事件订阅}


\texttt{subscribeEmbeddedPiSession()} 订阅 pi 的 \texttt{AgentSession} 事件：

\begin{lstlisting}[language=typescript]
const subscription = subscribeEmbeddedPiSession({
  session: activeSession,
  runId: params.runId,
  verboseLevel: params.verboseLevel,
  reasoningMode: params.reasoningLevel,
  toolResultFormat: params.toolResultFormat,
  onToolResult: params.onToolResult,
  onReasoningStream: params.onReasoningStream,
  onBlockReply: params.onBlockReply,
  onPartialReply: params.onPartialReply,
  onAgentEvent: params.onAgentEvent,
});
\end{lstlisting}


处理的事件包括：

\begin{itemize}
  \item \texttt{message\_start} / \texttt{message\_end} / \texttt{message\_update}（流式文本/思考）
  \item \texttt{tool\_execution\_start} / \texttt{tool\_execution\_update} / \texttt{tool\_execution\_end}
  \item \texttt{turn\_start} / \texttt{turn\_end}
  \item \texttt{agent\_start} / \texttt{agent\_end}
  \item \texttt{auto\_compaction\_start} / \texttt{auto\_compaction\_end}
\end{itemize}


\subsubsection{4. 提示}


设置完成后，会话被提示：

\begin{lstlisting}[language=typescript]
await session.prompt(effectivePrompt, { images: imageResult.images });
\end{lstlisting}


SDK 处理完整的智能体循环：发送到 LLM、执行工具调用、流式响应。

\subsection{工具架构}


\subsubsection{工具管道}


\begin{enumerate}
  \item \textbf{基础工具}：pi 的 \texttt{codingTools}（read、bash、edit、write）
  \item \textbf{自定义替换}：OpenClaw 将 bash 替换为 \texttt{exec}/\texttt{process}，为沙箱自定义 read/edit/write
  \item \textbf{OpenClaw 工具}：消息、浏览器、画布、会话、定时任务、Gateway 网关等
  \item \textbf{渠道工具}：Discord/Telegram/Slack/WhatsApp 特定的操作工具
  \item \textbf{策略过滤}：工具按配置文件、提供商、智能体、群组、沙箱策略过滤
  \item \textbf{Schema 规范化}：为 Gemini/OpenAI 的特殊情况清理 Schema
  \item \textbf{AbortSignal 包装}：工具被包装以尊重中止信号
\end{enumerate}


\subsubsection{工具定义适配器}


pi-agent-core 的 \texttt{AgentTool} 与 pi-coding-agent 的 \texttt{ToolDefinition} 有不同的 \texttt{execute} 签名。\texttt{pi-tool-definition-adapter.ts} 中的适配器桥接了这一点：

\begin{lstlisting}[language=typescript]
export function toToolDefinitions(tools: AnyAgentTool[]): ToolDefinition[] {
  return tools.map((tool) => ({
    name: tool.name,
    label: tool.label ?? name,
    description: tool.description ?? "",
    parameters: tool.parameters,
    execute: async (toolCallId, params, onUpdate, _ctx, signal) => {
      // pi-coding-agent signature differs from pi-agent-core
      return await tool.execute(toolCallId, params, signal, onUpdate);
    },
  }));
}
\end{lstlisting}


\subsubsection{工具拆分策略}


\texttt{splitSdkTools()} 通过 \texttt{customTools} 传递所有工具：

\begin{lstlisting}[language=typescript]
export function splitSdkTools(options: { tools: AnyAgentTool[]; sandboxEnabled: boolean }) {
  return {
    builtInTools: [], // Empty. We override everything
    customTools: toToolDefinitions(options.tools),
  };
}
\end{lstlisting}


这确保 OpenClaw 的策略过滤、沙箱集成和扩展工具集在各提供商之间保持一致。

\subsection{系统提示构建}


系统提示在 \texttt{buildAgentSystemPrompt()}（\texttt{system-prompt.ts}）中构建。它组装一个完整的提示，包含工具、工具调用风格、安全护栏、OpenClaw CLI 参考、Skills、文档、工作区、沙箱、消息、回复标签、语音、静默回复、心跳、运行时元数据等部分，以及启用时的记忆和反应，还有可选的上下文文件和额外系统提示内容。部分内容在子智能体使用的最小提示模式下会被裁剪。

提示在会话创建后通过 \texttt{applySystemPromptOverrideToSession()} 应用：

\begin{lstlisting}[language=typescript]
const systemPromptOverride = createSystemPromptOverride(appendPrompt);
applySystemPromptOverrideToSession(session, systemPromptOverride);
\end{lstlisting}


\subsection{会话管理}


\subsubsection{会话文件}


会话是具有树结构（id/parentId 链接）的 JSONL 文件。Pi 的 \texttt{SessionManager} 处理持久化：

\begin{lstlisting}[language=typescript]
const sessionManager = SessionManager.open(params.sessionFile);
\end{lstlisting}


OpenClaw 用 \texttt{guardSessionManager()} 包装它以确保工具结果安全。

\subsubsection{会话缓存}


\texttt{session-manager-cache.ts} 缓存 SessionManager 实例以避免重复的文件解析：

\begin{lstlisting}[language=typescript]
await prewarmSessionFile(params.sessionFile);
sessionManager = SessionManager.open(params.sessionFile);
trackSessionManagerAccess(params.sessionFile);
\end{lstlisting}


\subsubsection{历史限制}


\texttt{limitHistoryTurns()} 根据渠道类型（私信 vs 群组）裁剪对话历史。

\subsubsection{压缩}


自动压缩在上下文溢出时触发。\texttt{compactEmbeddedPiSessionDirect()} 处理手动压缩：

\begin{lstlisting}[language=typescript]
const compactResult = await compactEmbeddedPiSessionDirect({
  sessionId, sessionFile, provider, model, ...
});
\end{lstlisting}


\subsection{认证与模型解析}


\subsubsection{认证配置文件}


OpenClaw 维护一个认证配置文件存储，每个提供商有多个 API 密钥：

\begin{lstlisting}[language=typescript]
const authStore = ensureAuthProfileStore(agentDir, { allowKeychainPrompt: false });
const profileOrder = resolveAuthProfileOrder({ cfg, store: authStore, provider, preferredProfile });
\end{lstlisting}


配置文件在失败时轮换，并带有冷却跟踪：

\begin{lstlisting}[language=typescript]
await markAuthProfileFailure({ store, profileId, reason, cfg, agentDir });
const rotated = await advanceAuthProfile();
\end{lstlisting}


\subsubsection{模型解析}


\begin{lstlisting}[language=typescript]
import { resolveModel } from "./pi-embedded-runner/model.js";

const { model, error, authStorage, modelRegistry } = resolveModel(
  provider,
  modelId,
  agentDir,
  config,
);

// Uses pi's ModelRegistry and AuthStorage
authStorage.setRuntimeApiKey(model.provider, apiKeyInfo.apiKey);
\end{lstlisting}


\subsubsection{故障转移}


\texttt{FailoverError} 在配置了回退时触发模型回退：

\begin{lstlisting}[language=typescript]
if (fallbackConfigured && isFailoverErrorMessage(errorText)) {
  throw new FailoverError(errorText, {
    reason: promptFailoverReason ?? "unknown",
    provider,
    model: modelId,
    profileId,
    status: resolveFailoverStatus(promptFailoverReason),
  });
}
\end{lstlisting}


\subsection{Pi 扩展}


OpenClaw 加载自定义 pi 扩展以实现特殊行为：

\subsubsection{压缩安全护栏}


\texttt{pi-extensions/compaction-safeguard.ts} 为压缩添加护栏，包括自适应令牌预算以及工具失败和文件操作摘要：

\begin{lstlisting}[language=typescript]
if (resolveCompactionMode(params.cfg) === "safeguard") {
  setCompactionSafeguardRuntime(params.sessionManager, { maxHistoryShare });
  paths.push(resolvePiExtensionPath("compaction-safeguard"));
}
\end{lstlisting}


\subsubsection{上下文裁剪}


\texttt{pi-extensions/context-pruning.ts} 实现基于缓存 TTL 的上下文裁剪：

\begin{lstlisting}[language=typescript]
if (cfg?.agents?.defaults?.contextPruning?.mode === "cache-ttl") {
  setContextPruningRuntime(params.sessionManager, {
    settings,
    contextWindowTokens,
    isToolPrunable,
    lastCacheTouchAt,
  });
  paths.push(resolvePiExtensionPath("context-pruning"));
}
\end{lstlisting}


\subsection{流式传输与块回复}


\subsubsection{块分块}


\texttt{EmbeddedBlockChunker} 管理将流式文本分成离散的回复块：

\begin{lstlisting}[language=typescript]
const blockChunker = blockChunking ? new EmbeddedBlockChunker(blockChunking) : null;
\end{lstlisting}


\subsubsection{思考/最终标签剥离}


流式输出被处理以剥离 `\texttt{/}\texttt{ 块并提取 }` 内容：

\begin{lstlisting}[language=typescript]
const stripBlockTags = (text: string, state: { thinking: boolean; final: boolean }) => {
  // Strip <think>...</think> content
  // If enforceFinalTag, only return <final>...</final> content
};
\end{lstlisting}


\subsubsection{回复指令}


回复指令如 \texttt{[[media:url]]}、\texttt{[[voice]]}、\texttt{[[reply:id]]} 被解析和提取：

\begin{lstlisting}[language=typescript]
const { text: cleanedText, mediaUrls, audioAsVoice, replyToId } = consumeReplyDirectives(chunk);
\end{lstlisting}


\subsection{错误处理}


\subsubsection{错误分类}


\texttt{pi-embedded-helpers.ts} 对错误进行分类以进行适当处理：

\begin{lstlisting}[language=typescript]
isContextOverflowError(errorText)     // Context too large
isCompactionFailureError(errorText)   // Compaction failed
isAuthAssistantError(lastAssistant)   // Auth failure
isRateLimitAssistantError(...)        // Rate limited
isFailoverAssistantError(...)         // Should failover
classifyFailoverReason(errorText)     // "auth" | "rate_limit" | "quota" | "timeout" | ...
\end{lstlisting}


\subsubsection{思考级别回退}


如果思考级别不受支持，它会回退：

\begin{lstlisting}[language=typescript]
const fallbackThinking = pickFallbackThinkingLevel({
  message: errorText,
  attempted: attemptedThinking,
});
if (fallbackThinking) {
  thinkLevel = fallbackThinking;
  continue;
}
\end{lstlisting}


\subsection{沙箱集成}


当启用沙箱模式时，工具和路径受到约束：

\begin{lstlisting}[language=typescript]
const sandbox = await resolveSandboxContext({
  config: params.config,
  sessionKey: sandboxSessionKey,
  workspaceDir: resolvedWorkspace,
});

if (sandboxRoot) {
  // Use sandboxed read/edit/write tools
  // Exec runs in container
  // Browser uses bridge URL
}
\end{lstlisting}


\subsection{提供商特定处理}


\subsubsection{Anthropic}


\begin{itemize}
  \item 拒绝魔术字符串清除
  \item 连续角色的回合验证
  \item Claude Code 参数兼容性
\end{itemize}


\subsubsection{Google/Gemini}


\begin{itemize}
  \item 回合排序修复（\texttt{applyGoogleTurnOrderingFix}）
  \item 工具 schema 清理（\texttt{sanitizeToolsForGoogle}）
  \item 会话历史清理（\texttt{sanitizeSessionHistory}）
\end{itemize}


\subsubsection{OpenAI}


\begin{itemize}
  \item Codex 模型的 \texttt{apply\_patch} 工具
  \item 思考级别降级处理
\end{itemize}


\subsection{TUI 集成}


OpenClaw 还有一个本地 TUI 模式，直接使用 pi-tui 组件：

\begin{lstlisting}[language=typescript]
// src/tui/tui.ts
import { ... } from "@mariozechner/pi-tui";
\end{lstlisting}


这提供了与 pi 原生模式类似的交互式终端体验。

\subsection{与 Pi CLI 的主要区别}



\begin{table}[h]
\centering
\small
\begin{tabular}{|l|l|l|}
\hline
方面 & Pi CLI & OpenClaw 嵌入式 \\
\hline
调用方式 & \texttt{pi} 命令 / RPC & 通过 \texttt{createAgentSession()} 的 SDK \\
\hline
工具 & 默认编码工具 & 自定义 OpenClaw 工具套件 \\
\hline
系统提示 & AGENTS.md + prompts & 按渠道/上下文动态生成 \\
\hline
会话存储 & \texttt{\textasciitilde{}/.pi/agent/sessions/} & \texttt{\textasciitilde{}/.openclaw/agents//sessions/}（或 \texttt{\$OPENCLAW\_STATE\_DIR/agents//sessions/}） \\
\hline
认证 & 单一凭证 & 带轮换的多配置文件 \\
\hline
扩展 & 从磁盘加载 & 编程方式 + 磁盘路径 \\
\hline
事件处理 & TUI 渲染 & 基于回调（onBlockReply 等） \\
\hline
\end{tabular}
\end{table}



\subsection{未来考虑}


可能需要重构的领域：

\begin{enumerate}
  \item \textbf{工具签名对齐}：目前在 pi-agent-core 和 pi-coding-agent 签名之间适配
  \item \textbf{会话管理器包装}：\texttt{guardSessionManager} 增加了安全性但增加了复杂性
  \item \textbf{扩展加载}：可以更直接地使用 pi 的 \texttt{ResourceLoader}
  \item \textbf{流式处理器复杂性}：\texttt{subscribeEmbeddedPiSession} 已经变得很大
  \item \textbf{提供商特殊情况}：许多提供商特定的代码路径，pi 可能可以处理
\end{enumerate}


\subsection{测试}


所有涵盖 pi 集成及其扩展的现有测试：

\begin{itemize}
  \item \texttt{src/agents/pi-embedded-block-chunker.test.ts}
  \item \texttt{src/agents/pi-embedded-helpers.buildbootstrapcontextfiles.test.ts}
  \item \texttt{src/agents/pi-embedded-helpers.classifyfailoverreason.test.ts}
  \item \texttt{src/agents/pi-embedded-helpers.downgradeopenai-reasoning.test.ts}
  \item \texttt{src/agents/pi-embedded-helpers.formatassistanterrortext.test.ts}
  \item \texttt{src/agents/pi-embedded-helpers.formatrawassistanterrorforui.test.ts}
  \item \texttt{src/agents/pi-embedded-helpers.image-dimension-error.test.ts}
  \item \texttt{src/agents/pi-embedded-helpers.image-size-error.test.ts}
  \item \texttt{src/agents/pi-embedded-helpers.isautherrormessage.test.ts}
  \item \texttt{src/agents/pi-embedded-helpers.isbillingerrormessage.test.ts}
  \item \texttt{src/agents/pi-embedded-helpers.iscloudcodeassistformaterror.test.ts}
  \item \texttt{src/agents/pi-embedded-helpers.iscompactionfailureerror.test.ts}
  \item \texttt{src/agents/pi-embedded-helpers.iscontextoverflowerror.test.ts}
  \item \texttt{src/agents/pi-embedded-helpers.isfailovererrormessage.test.ts}
  \item \texttt{src/agents/pi-embedded-helpers.islikelycontextoverflowerror.test.ts}
  \item \texttt{src/agents/pi-embedded-helpers.ismessagingtoolduplicate.test.ts}
  \item \texttt{src/agents/pi-embedded-helpers.messaging-duplicate.test.ts}
  \item \texttt{src/agents/pi-embedded-helpers.normalizetextforcomparison.test.ts}
  \item \texttt{src/agents/pi-embedded-helpers.resolvebootstrapmaxchars.test.ts}
  \item \texttt{src/agents/pi-embedded-helpers.sanitize-session-messages-images.keeps-tool-call-tool-result-ids-unchanged.test.ts}
  \item \texttt{src/agents/pi-embedded-helpers.sanitize-session-messages-images.removes-empty-assistant-text-blocks-but-preserves.test.ts}
  \item \texttt{src/agents/pi-embedded-helpers.sanitizegoogleturnordering.test.ts}
  \item \texttt{src/agents/pi-embedded-helpers.sanitizesessionmessagesimages-thought-signature-stripping.test.ts}
  \item \texttt{src/agents/pi-embedded-helpers.sanitizetoolcallid.test.ts}
  \item \texttt{src/agents/pi-embedded-helpers.sanitizeuserfacingtext.test.ts}
  \item \texttt{src/agents/pi-embedded-helpers.stripthoughtsignatures.test.ts}
  \item \texttt{src/agents/pi-embedded-helpers.validate-turns.test.ts}
  \item \texttt{src/agents/pi-embedded-runner-extraparams.live.test.ts}（实时）
  \item \texttt{src/agents/pi-embedded-runner-extraparams.test.ts}
  \item \texttt{src/agents/pi-embedded-runner.applygoogleturnorderingfix.test.ts}
  \item \texttt{src/agents/pi-embedded-runner.buildembeddedsandboxinfo.test.ts}
  \item \texttt{src/agents/pi-embedded-runner.createsystempromptoverride.test.ts}
  \item \texttt{src/agents/pi-embedded-runner.get-dm-history-limit-from-session-key.falls-back-provider-default-per-dm-not.test.ts}
  \item \texttt{src/agents/pi-embedded-runner.get-dm-history-limit-from-session-key.returns-undefined-sessionkey-is-undefined.test.ts}
  \item \texttt{src/agents/pi-embedded-runner.google-sanitize-thinking.test.ts}
  \item \texttt{src/agents/pi-embedded-runner.guard.test.ts}
  \item \texttt{src/agents/pi-embedded-runner.limithistoryturns.test.ts}
  \item \texttt{src/agents/pi-embedded-runner.resolvesessionagentids.test.ts}
  \item \texttt{src/agents/pi-embedded-runner.run-embedded-pi-agent.auth-profile-rotation.test.ts}
  \item \texttt{src/agents/pi-embedded-runner.sanitize-session-history.test.ts}
  \item \texttt{src/agents/pi-embedded-runner.splitsdktools.test.ts}
  \item \texttt{src/agents/pi-embedded-runner.test.ts}
  \item \texttt{src/agents/pi-embedded-subscribe.code-span-awareness.test.ts}
  \item \texttt{src/agents/pi-embedded-subscribe.reply-tags.test.ts}
  \item \texttt{src/agents/pi-embedded-subscribe.subscribe-embedded-pi-session.calls-onblockreplyflush-before-tool-execution-start-preserve.test.ts}
  \item \texttt{src/agents/pi-embedded-subscribe.subscribe-embedded-pi-session.does-not-append-text-end-content-is.test.ts}
  \item \texttt{src/agents/pi-embedded-subscribe.subscribe-embedded-pi-session.does-not-call-onblockreplyflush-callback-is-not.test.ts}
  \item \texttt{src/agents/pi-embedded-subscribe.subscribe-embedded-pi-session.does-not-duplicate-text-end-repeats-full.test.ts}
  \item \texttt{src/agents/pi-embedded-subscribe.subscribe-embedded-pi-session.does-not-emit-duplicate-block-replies-text.test.ts}
  \item \texttt{src/agents/pi-embedded-subscribe.subscribe-embedded-pi-session.emits-block-replies-text-end-does-not.test.ts}
  \item \texttt{src/agents/pi-embedded-subscribe.subscribe-embedded-pi-session.emits-reasoning-as-separate-message-enabled.test.ts}
  \item \texttt{src/agents/pi-embedded-subscribe.subscribe-embedded-pi-session.filters-final-suppresses-output-without-start-tag.test.ts}
  \item \texttt{src/agents/pi-embedded-subscribe.subscribe-embedded-pi-session.includes-canvas-action-metadata-tool-summaries.test.ts}
  \item \texttt{src/agents/pi-embedded-subscribe.subscribe-embedded-pi-session.keeps-assistanttexts-final-answer-block-replies-are.test.ts}
  \item \texttt{src/agents/pi-embedded-subscribe.subscribe-embedded-pi-session.keeps-indented-fenced-blocks-intact.test.ts}
  \item \texttt{src/agents/pi-embedded-subscribe.subscribe-embedded-pi-session.reopens-fenced-blocks-splitting-inside-them.test.ts}
  \item \texttt{src/agents/pi-embedded-subscribe.subscribe-embedded-pi-session.splits-long-single-line-fenced-blocks-reopen.test.ts}
  \item \texttt{src/agents/pi-embedded-subscribe.subscribe-embedded-pi-session.streams-soft-chunks-paragraph-preference.test.ts}
  \item \texttt{src/agents/pi-embedded-subscribe.subscribe-embedded-pi-session.subscribeembeddedpisession.test.ts}
  \item \texttt{src/agents/pi-embedded-subscribe.subscribe-embedded-pi-session.suppresses-message-end-block-replies-message-tool.test.ts}
  \item \texttt{src/agents/pi-embedded-subscribe.subscribe-embedded-pi-session.waits-multiple-compaction-retries-before-resolving.test.ts}
  \item \texttt{src/agents/pi-embedded-subscribe.tools.test.ts}
  \item \texttt{src/agents/pi-embedded-utils.test.ts}
  \item \texttt{src/agents/pi-extensions/compaction-safeguard.test.ts}
  \item \texttt{src/agents/pi-extensions/context-pruning.test.ts}
  \item \texttt{src/agents/pi-settings.test.ts}
  \item \texttt{src/agents/pi-tool-definition-adapter.test.ts}
  \item \texttt{src/agents/pi-tools-agent-config.test.ts}
  \item \texttt{src/agents/pi-tools.create-openclaw-coding-tools.adds-claude-style-aliases-schemas-without-dropping-b.test.ts}
  \item \texttt{src/agents/pi-tools.create-openclaw-coding-tools.adds-claude-style-aliases-schemas-without-dropping-d.test.ts}
  \item \texttt{src/agents/pi-tools.create-openclaw-coding-tools.adds-claude-style-aliases-schemas-without-dropping-f.test.ts}
  \item \texttt{src/agents/pi-tools.create-openclaw-coding-tools.adds-claude-style-aliases-schemas-without-dropping.test.ts}
  \item \texttt{src/agents/pi-tools.policy.test.ts}
  \item \texttt{src/agents/pi-tools.safe-bins.test.ts}
  \item \texttt{src/agents/pi-tools.workspace-paths.test.ts}
\end{itemize}



\clearpage

% === 源文件: prose.md ===
\section{OpenProse}


OpenProse 是一种可移植的、以 Markdown 为中心的工作流格式，用于编排 AI 会话。在 OpenClaw 中，它作为插件发布，安装一个 OpenProse Skills 包以及一个 \texttt{/prose} 斜杠命令。程序存放在 \texttt{.prose} 文件中，可以生成多个具有显式控制流的子智能体。

官方网站：https://www.prose.md

\subsection{它能做什么}


\begin{itemize}
  \item 具有显式并行性的多智能体研究 + 综合。
  \item 可重复的批准安全工作流（代码审查、事件分类、内容管道）。
  \item 可在支持的智能体运行时之间运行的可重用 \texttt{.prose} 程序。
\end{itemize}


\subsection{安装 + 启用}


捆绑的插件默认是禁用的。启用 OpenProse：

\begin{lstlisting}[language=bash]
openclaw plugins enable open-prose
\end{lstlisting}


启用插件后重启 Gateway 网关。

开发/本地检出：\texttt{openclaw plugins install ./extensions/open-prose}

相关文档：\href{/tools/plugin}{插件}、\href{/plugins/manifest}{插件清单}、\href{/tools/skills}{Skills}。

\subsection{斜杠命令}


OpenProse 将 \texttt{/prose} 注册为用户可调用的 Skills 命令。它路由到 OpenProse VM 指令，并在底层使用 OpenClaw 工具。

常用命令：

\begin{lstlisting}
/prose help
/prose run <file.prose>
/prose run <handle/slug>
/prose run <https://example.com/file.prose>
/prose compile <file.prose>
/prose examples
/prose update
\end{lstlisting}


\subsection{示例：一个简单的 .prose 文件}


\begin{lstlisting}
# Research + synthesis with two agents running in parallel.

input topic: "What should we research?"

agent researcher:
  model: sonnet
  prompt: "You research thoroughly and cite sources."

agent writer:
  model: opus
  prompt: "You write a concise summary."

parallel:
  findings = session: researcher
    prompt: "Research {topic}."
  draft = session: writer
    prompt: "Summarize {topic}."

session "Merge the findings + draft into a final answer."
context: { findings, draft }
\end{lstlisting}


\subsection{文件位置}


OpenProse 将状态保存在工作空间的 \texttt{.prose/} 下：

\begin{lstlisting}
.prose/
├── .env
├── runs/
│   └── {YYYYMMDD}-{HHMMSS}-{random}/
│       ├── program.prose
│       ├── state.md
│       ├── bindings/
│       └── agents/
└── agents/
\end{lstlisting}


用户级持久智能体位于：

\begin{lstlisting}
~/.prose/agents/
\end{lstlisting}


\subsection{状态模式}


OpenProse 支持多种状态后端：

\begin{itemize}
  \item \textbf{filesystem}（默认）：\texttt{.prose/runs/...}
  \item \textbf{in-context}：瞬态，用于小程序
  \item \textbf{sqlite}（实验性）：需要 \texttt{sqlite3} 二进制文件
  \item \textbf{postgres}（实验性）：需要 \texttt{psql} 和连接字符串
\end{itemize}


说明：

\begin{itemize}
  \item sqlite/postgres 是选择加入的，且处于实验阶段。
  \item postgres 凭证会流入子智能体日志；请使用专用的、最小权限的数据库。
\end{itemize}


\subsection{远程程序}


\texttt{/prose run } 解析为 \texttt{https://p.prose.md//}。
直接 URL 按原样获取。这使用 \texttt{web\_fetch} 工具（或用于 POST 的 \texttt{exec}）。

\subsection{OpenClaw 运行时映射}


OpenProse 程序映射到 OpenClaw 原语：


\begin{table}[h]
\centering
\small
\begin{tabular}{|l|l|}
\hline
OpenProse 概念 & OpenClaw 工具 \\
\hline
生成会话 / Task 工具 & \texttt{sessions\_spawn} \\
\hline
文件读/写 & \texttt{read} / \texttt{write} \\
\hline
Web 获取 & \texttt{web\_fetch} \\
\hline
\end{tabular}
\end{table}



如果你的工具白名单阻止这些工具，OpenProse 程序将失败。参见 \href{/tools/skills-config}{Skills 配置}。

\subsection{安全 + 批准}


将 \texttt{.prose} 文件视为代码。运行前请审查。使用 OpenClaw 工具白名单和批准门控来控制副作用。

对于确定性的、批准门控的工作流，可与 \href{/tools/lobster}{Lobster} 比较。


\clearpage

% === 源文件: tts.md ===
\section{文本转语音（TTS）}


OpenClaw 可以使用 ElevenLabs、OpenAI 或 Edge TTS 将出站回复转换为音频。它可以在任何 OpenClaw 能发送音频的地方工作；Telegram 会显示圆形语音消息气泡。

\subsection{支持的服务}


\begin{itemize}
  \item \textbf{ElevenLabs}（主要或备用提供商）
  \item \textbf{OpenAI}（主要或备用提供商；也用于摘要）
  \item \textbf{Edge TTS}（主要或备用提供商；使用 \texttt{node-edge-tts}，无 API 密钥时为默认）
\end{itemize}


\subsubsection{Edge TTS 注意事项}


Edge TTS 通过 \texttt{node-edge-tts} 库使用 Microsoft Edge 的在线神经网络 TTS 服务。它是托管服务（非本地），使用 Microsoft 的端点，不需要 API 密钥。\texttt{node-edge-tts} 公开了语音配置选项和输出格式，但并非所有选项都被 Edge 服务支持。citeturn2search0

由于 Edge TTS 是一个没有公布 SLA 或配额的公共 Web 服务，请将其视为尽力而为。如果你需要有保证的限制和支持，请使用 OpenAI 或 ElevenLabs。Microsoft 的语音 REST API 记录了每个请求 10 分钟的音频限制；Edge TTS 没有公布限制，所以假设类似或更低的限制。citeturn0search3

\subsection{可选密钥}


如果你想使用 OpenAI 或 ElevenLabs：

\begin{itemize}
  \item \texttt{ELEVENLABS\_API\_KEY}（或 \texttt{XI\_API\_KEY}）
  \item \texttt{OPENAI\_API\_KEY}
\end{itemize}


Edge TTS \textbf{不}需要 API 密钥。如果没有找到 API 密钥，OpenClaw 默认使用 Edge TTS（除非通过 \texttt{messages.tts.edge.enabled=false} 禁用）。

如果配置了多个提供商，首先使用选定的提供商，其他作为备用选项。自动摘要使用配置的 \texttt{summaryModel}（或 \texttt{agents.defaults.model.primary}），所以如果你启用摘要，该提供商也必须经过认证。

\subsection{服务链接}


\begin{itemize}
  \item \href{https://platform.openai.com/docs/guides/text-to-speech}{OpenAI 文本转语音指南}
  \item \href{https://platform.openai.com/docs/api-reference/audio}{OpenAI 音频 API 参考}
  \item \href{https://elevenlabs.io/docs/api-reference/text-to-speech}{ElevenLabs 文本转语音}
  \item \href{https://elevenlabs.io/docs/api-reference/authentication}{ElevenLabs 认证}
  \item \href{https://github.com/SchneeHertz/node-edge-tts}{node-edge-tts}
  \item \href{https://learn.microsoft.com/azure/ai-services/speech-service/rest-text-to-speech\#audio-outputs}{Microsoft 语音输出格式}
\end{itemize}


\subsection{默认启用吗？}


不是。自动 TTS 默认\textbf{关闭}。在配置中使用 \texttt{messages.tts.auto} 或在每个会话中使用 \texttt{/tts always}（别名：\texttt{/tts on}）启用它。

一旦 TTS 开启，Edge TTS \textbf{是}默认启用的，并在没有 OpenAI 或 ElevenLabs API 密钥时自动使用。

\subsection{配置}


TTS 配置位于 \texttt{openclaw.json} 中的 \texttt{messages.tts} 下。完整 schema 在 \href{/gateway/configuration}{Gateway 网关配置}中。

\subsubsection{最小配置（启用 + 提供商）}


\begin{lstlisting}[language=json]
{
  messages: {
    tts: {
      auto: "always",
      provider: "elevenlabs",
    },
  },
}
\end{lstlisting}


\subsubsection{OpenAI 主要，ElevenLabs 备用}


\begin{lstlisting}[language=json]
{
  messages: {
    tts: {
      auto: "always",
      provider: "openai",
      summaryModel: "openai/gpt-4.1-mini",
      modelOverrides: {
        enabled: true,
      },
      openai: {
        apiKey: "openai_api_key",
        model: "gpt-4o-mini-tts",
        voice: "alloy",
      },
      elevenlabs: {
        apiKey: "elevenlabs_api_key",
        baseUrl: "https://api.elevenlabs.io",
        voiceId: "voice_id",
        modelId: "eleven_multilingual_v2",
        seed: 42,
        applyTextNormalization: "auto",
        languageCode: "en",
        voiceSettings: {
          stability: 0.5,
          similarityBoost: 0.75,
          style: 0.0,
          useSpeakerBoost: true,
          speed: 1.0,
        },
      },
    },
  },
}
\end{lstlisting}


\subsubsection{Edge TTS 主要（无 API 密钥）}


\begin{lstlisting}[language=json]
{
  messages: {
    tts: {
      auto: "always",
      provider: "edge",
      edge: {
        enabled: true,
        voice: "en-US-MichelleNeural",
        lang: "en-US",
        outputFormat: "audio-24khz-48kbitrate-mono-mp3",
        rate: "+10%",
        pitch: "-5%",
      },
    },
  },
}
\end{lstlisting}


\subsubsection{禁用 Edge TTS}


\begin{lstlisting}[language=json]
{
  messages: {
    tts: {
      edge: {
        enabled: false,
      },
    },
  },
}
\end{lstlisting}


\subsubsection{自定义限制 + 偏好路径}


\begin{lstlisting}[language=json]
{
  messages: {
    tts: {
      auto: "always",
      maxTextLength: 4000,
      timeoutMs: 30000,
      prefsPath: "~/.openclaw/settings/tts.json",
    },
  },
}
\end{lstlisting}


\subsubsection{仅在收到语音消息后用音频回复}


\begin{lstlisting}[language=json]
{
  messages: {
    tts: {
      auto: "inbound",
    },
  },
}
\end{lstlisting}


\subsubsection{禁用长回复的自动摘要}


\begin{lstlisting}[language=json]
{
  messages: {
    tts: {
      auto: "always",
    },
  },
}
\end{lstlisting}


然后运行：

\begin{lstlisting}
/tts summary off
\end{lstlisting}


\subsubsection{字段说明}


\begin{itemize}
  \item \texttt{auto}：自动 TTS 模式（\texttt{off}、\texttt{always}、\texttt{inbound}、\texttt{tagged}）。
  \item \texttt{inbound} 仅在收到语音消息后发送音频。
  \item \texttt{tagged} 仅在回复包含 \texttt{[[tts]]} 标签时发送音频。
  \item \texttt{enabled}：旧版开关（doctor 将其迁移到 \texttt{auto}）。
  \item \texttt{mode}：\texttt{"final"}（默认）或 \texttt{"all"}（包括工具/分块回复）。
  \item \texttt{provider}：\texttt{"elevenlabs"}、\texttt{"openai"} 或 \texttt{"edge"}（自动备用）。
  \item 如果 \texttt{provider} \textbf{未设置}，OpenClaw 优先选择 \texttt{openai}（如果有密钥），然后是 \texttt{elevenlabs}（如果有密钥），否则是 \texttt{edge}。
  \item \texttt{summaryModel}：用于自动摘要的可选廉价模型；默认为 \texttt{agents.defaults.model.primary}。
  \item 接受 \texttt{provider/model} 或配置的模型别名。
  \item \texttt{modelOverrides}：允许模型发出 TTS 指令（默认开启）。
  \item \texttt{maxTextLength}：TTS 输入的硬性上限（字符）。超出时 \texttt{/tts audio} 会失败。
  \item \texttt{timeoutMs}：请求超时（毫秒）。
  \item \texttt{prefsPath}：覆盖本地偏好 JSON 路径（提供商/限制/摘要）。
  \item \texttt{apiKey} 值回退到环境变量（\texttt{ELEVENLABS\_API\_KEY}/\texttt{XI\_API\_KEY}、\texttt{OPENAI\_API\_KEY}）。
  \item \texttt{elevenlabs.baseUrl}：覆盖 ElevenLabs API 基础 URL。
  \item \texttt{elevenlabs.voiceSettings}：
  \item \texttt{stability}、\texttt{similarityBoost}、\texttt{style}：\texttt{0..1}
  \item \texttt{useSpeakerBoost}：\texttt{true|false}
  \item \texttt{speed}：\texttt{0.5..2.0}（1.0 = 正常）
  \item \texttt{elevenlabs.applyTextNormalization}：\texttt{auto|on|off}
  \item \texttt{elevenlabs.languageCode}：2 字母 ISO 639-1（例如 \texttt{en}、\texttt{de}）
  \item \texttt{elevenlabs.seed}：整数 \texttt{0..4294967295}（尽力确定性）
  \item \texttt{edge.enabled}：允许 Edge TTS 使用（默认 \texttt{true}；无 API 密钥）。
  \item \texttt{edge.voice}：Edge 神经网络语音名称（例如 \texttt{en-US-MichelleNeural}）。
  \item \texttt{edge.lang}：语言代码（例如 \texttt{en-US}）。
  \item \texttt{edge.outputFormat}：Edge 输出格式（例如 \texttt{audio-24khz-48kbitrate-mono-mp3}）。
  \item 有效值参见 Microsoft 语音输出格式；并非所有格式都被 Edge 支持。
  \item \texttt{edge.rate} / \texttt{edge.pitch} / \texttt{edge.volume}：百分比字符串（例如 \texttt{+10\%}、\texttt{-5\%}）。
  \item \texttt{edge.saveSubtitles}：在音频文件旁边写入 JSON 字幕。
  \item \texttt{edge.proxy}：Edge TTS 请求的代理 URL。
  \item \texttt{edge.timeoutMs}：请求超时覆盖（毫秒）。
\end{itemize}


\subsection{模型驱动覆盖（默认开启）}


默认情况下，模型\textbf{可以}为单个回复发出 TTS 指令。当 \texttt{messages.tts.auto} 为 \texttt{tagged} 时，需要这些指令来触发音频。

启用后，模型可以发出 \texttt{[[tts:...]]} 指令来覆盖单个回复的语音，加上可选的 \texttt{[[tts:text]]...[[/tts:text]]} 块来提供表达性标签（笑声、唱歌提示等），这些仅应出现在音频中。

示例回复负载：

\begin{lstlisting}
Here you go.

[[tts:provider=elevenlabs voiceId=pMsXgVXv3BLzUgSXRplE model=eleven_v3 speed=1.1]]
[[tts:text]](laughs) Read the song once more.[[/tts:text]]
\end{lstlisting}


可用指令键（启用时）：

\begin{itemize}
  \item \texttt{provider}（\texttt{openai} | \texttt{elevenlabs} | \texttt{edge}）
  \item \texttt{voice}（OpenAI 语音）或 \texttt{voiceId}（ElevenLabs）
  \item \texttt{model}（OpenAI TTS 模型或 ElevenLabs 模型 ID）
  \item \texttt{stability}、\texttt{similarityBoost}、\texttt{style}、\texttt{speed}、\texttt{useSpeakerBoost}
  \item \texttt{applyTextNormalization}（\texttt{auto|on|off}）
  \item \texttt{languageCode}（ISO 639-1）
  \item \texttt{seed}
\end{itemize}


禁用所有模型覆盖：

\begin{lstlisting}[language=json]
{
  messages: {
    tts: {
      modelOverrides: {
        enabled: false,
      },
    },
  },
}
\end{lstlisting}


可选白名单（禁用特定覆盖同时保持标签启用）：

\begin{lstlisting}[language=json]
{
  messages: {
    tts: {
      modelOverrides: {
        enabled: true,
        allowProvider: false,
        allowSeed: false,
      },
    },
  },
}
\end{lstlisting}


\subsection{单用户偏好}


斜杠命令将本地覆盖写入 \texttt{prefsPath}（默认：\texttt{\textasciitilde{}/.openclaw/settings/tts.json}，可通过 \texttt{OPENCLAW\_TTS\_PREFS} 或 \texttt{messages.tts.prefsPath} 覆盖）。

存储的字段：

\begin{itemize}
  \item \texttt{enabled}
  \item \texttt{provider}
  \item \texttt{maxLength}（摘要阈值；默认 1500 字符）
  \item \texttt{summarize}（默认 \texttt{true}）
\end{itemize}


这些为该主机覆盖 \texttt{messages.tts.*}。

\subsection{输出格式（固定）}


\begin{itemize}
  \item \textbf{Telegram}：Opus 语音消息（ElevenLabs 的 \texttt{opus\_48000\_64}，OpenAI 的 \texttt{opus}）。
  \item 48kHz / 64kbps 是语音消息的良好权衡，圆形气泡所必需。
  \item \textbf{其他渠道}：MP3（ElevenLabs 的 \texttt{mp3\_44100\_128}，OpenAI 的 \texttt{mp3}）。
  \item 44.1kHz / 128kbps 是语音清晰度的默认平衡。
  \item \textbf{Edge TTS}：使用 \texttt{edge.outputFormat}（默认 \texttt{audio-24khz-48kbitrate-mono-mp3}）。
  \item \texttt{node-edge-tts} 接受 \texttt{outputFormat}，但并非所有格式都可从 Edge 服务获得。citeturn2search0
  \item 输出格式值遵循 Microsoft 语音输出格式（包括 Ogg/WebM Opus）。citeturn1search0
  \item Telegram \texttt{sendVoice} 接受 OGG/MP3/M4A；如果你需要有保证的 Opus 语音消息，请使用 OpenAI/ElevenLabs。citeturn1search1
  \item 如果配置的 Edge 输出格式失败，OpenClaw 会使用 MP3 重试。
\end{itemize}


OpenAI/ElevenLabs 格式是固定的；Telegram 期望 Opus 以获得语音消息用户体验。

\subsection{自动 TTS 行为}


启用后，OpenClaw：

\begin{itemize}
  \item 如果回复已包含媒体或 \texttt{MEDIA:} 指令，则跳过 TTS。
  \item 跳过非常短的回复（< 10 字符）。
  \item 启用时使用 \texttt{agents.defaults.model.primary}（或 \texttt{summaryModel}）对长回复进行摘要。
  \item 将生成的音频附加到回复中。
\end{itemize}


如果回复超过 \texttt{maxLength} 且摘要关闭（或没有摘要模型的 API 密钥），则跳过音频并发送正常的文本回复。

\subsection{流程图}


\begin{lstlisting}
回复 -> TTS 启用？
  否  -> 发送文本
  是  -> 有媒体 / MEDIA: / 太短？
          是 -> 发送文本
          否 -> 长度 > 限制？
                   否  -> TTS -> 附加音频
                   是  -> 摘要启用？
                            否  -> 发送文本
                            是  -> 摘要（summaryModel 或 agents.defaults.model.primary）
                                      -> TTS -> 附加音频
\end{lstlisting}


\subsection{斜杠命令用法}


只有一个命令：\texttt{/tts}。参见\href{/tools/slash-commands}{斜杠命令}了解启用详情。

Discord 注意：\texttt{/tts} 是 Discord 的内置命令，所以 OpenClaw 在那里注册 \texttt{/voice} 作为原生命令。文本 \texttt{/tts ...} 仍然有效。

\begin{lstlisting}
/tts off
/tts always
/tts inbound
/tts tagged
/tts status
/tts provider openai
/tts limit 2000
/tts summary off
/tts audio Hello from OpenClaw
\end{lstlisting}


注意事项：

\begin{itemize}
  \item 命令需要授权发送者（白名单/所有者规则仍然适用）。
  \item 必须启用 \texttt{commands.text} 或原生命令注册。
  \item \texttt{off|always|inbound|tagged} 是单会话开关（\texttt{/tts on} 是 \texttt{/tts always} 的别名）。
  \item \texttt{limit} 和 \texttt{summary} 存储在本地偏好中，不在主配置中。
  \item \texttt{/tts audio} 生成一次性音频回复（不会开启 TTS）。
\end{itemize}


\subsection{智能体工具}


\texttt{tts} 工具将文本转换为语音并返回 \texttt{MEDIA:} 路径。当结果与 Telegram 兼容时，工具包含 \texttt{[[audio\_as\_voice]]}，以便 Telegram 发送语音气泡。

\subsection{Gateway 网关 RPC}


Gateway 网关方法：

\begin{itemize}
  \item \texttt{tts.status}
  \item \texttt{tts.enable}
  \item \texttt{tts.disable}
  \item \texttt{tts.convert}
  \item \texttt{tts.setProvider}
  \item \texttt{tts.providers}
\end{itemize}



\clearpage

% === 源文件: vps.md ===
\section{VPS 托管}


本中心链接到支持的 VPS/托管指南，并在高层次上解释云部署的工作原理。

\subsection{选择提供商}


\begin{itemize}
  \item \textbf{Railway}（一键 + 浏览器设置）：\href{/install/railway}{Railway}
  \item \textbf{Northflank}（一键 + 浏览器设置）：\href{/install/northflank}{Northflank}
  \item \textbf{Oracle Cloud（永久免费）}：\href{/platforms/oracle}{Oracle} — \$0/月（永久免费，ARM；容量/注册可能不太稳定）
  \item \textbf{Fly.io}：\href{/install/fly}{Fly.io}
  \item \textbf{Hetzner（Docker）}：\href{/install/hetzner}{Hetzner}
  \item \textbf{GCP（Compute Engine）}：\href{/install/gcp}{GCP}
  \item \textbf{exe.dev}（VM + HTTPS 代理）：\href{/install/exe-dev}{exe.dev}
  \item \textbf{AWS（EC2/Lightsail/免费套餐）}：也运行良好。视频指南：
\end{itemize}

https://x.com/techfrenAJ/status/2014934471095812547

\subsection{云设置的工作原理}


\begin{itemize}
  \item \textbf{Gateway 网关运行在 VPS 上}并拥有状态 + 工作区。
  \item 你通过\textbf{控制 UI} 或 \textbf{Tailscale/SSH} 从笔记本电脑/手机连接。
  \item 将 VPS 视为数据源并\textbf{备份}状态 + 工作区。
  \item 安全默认：将 Gateway 网关保持在 loopback 上，通过 SSH 隧道或 Tailscale Serve 访问。
\end{itemize}

如果你绑定到 \texttt{lan}/\texttt{tailnet}，需要 \texttt{gateway.auth.token} 或 \texttt{gateway.auth.password}。

远程访问：\href{/gateway/remote}{Gateway 网关远程访问}
平台中心：\href{/platforms}{平台}

\subsection{在 VPS 上使用节点}


你可以将 Gateway 网关保持在云端，并在本地设备（Mac/iOS/Android/无头）上配对\textbf{节点}。节点提供本地屏幕/摄像头/canvas 和 \texttt{system.run} 功能，而 Gateway 网关保持在云端。

文档：\href{/nodes}{节点}，\href{/cli/nodes}{节点 CLI}


\clearpage
